\documentclass[10pt,a4paper]{amsart}
\usepackage[margin=19mm]{geometry}
\usepackage{amsmath,amssymb,mathtools}
\usepackage[scr=rsfs,cal=euler]{mathalfa}
\usepackage{xfrac}
\usepackage[unicode,hidelinks,bookmarks=false]{hyperref}
\newcommand{\ie}{i.e.\ }
\newcommand{\cf}{cf.\ }
\newcommand{\resp}{respectively\ }
\newcommand{\Subsection}{\subsection}
\providecommand{\nobreakdash}{\nobreak\hbox{-}}
\setlength{\emergencystretch}{2em}
\multlinegap0pt
\usepackage{stmaryrd}
\usepackage{amssymb}
\usepackage{float}
\usepackage{multirow}
\usepackage{tikz}
\usepackage{enumerate}
\usepackage{subfig}
\usepackage{algorithm}\usepackage{algpseudocode}
\newtheorem{lemma}{Lemma}
\newtheorem{theorem}{Theorem}
\newcommand\bE{\mathbf{E}}
\newcommand\bU{\mathbf{U}}
\newcommand\bV{\mathbf{V}}
\newcommand\bbeta{{\boldsymbol{\beta}}}
\newcommand\bx{\boldsymbol{x}}
\newcommand\calF{\mathcal{F}}
\newcommand\calP{{\mathcal{P}}}
\title{A finite element method preserving the eigenvalue range of symmetric tensor fields}
\author{Abdolreza Amiri, Gabriel R. Barrenechea, Tristan Pryer}
\date{Source: arXiv:2601.04839v1 (CC BY 4.0)}
\begin{document}
\maketitle

\noindent\emph{Source and licence.} A finite element method preserving the eigenvalue range of symmetric tensor fields. Version arXiv:2601.04839v1 (CC BY 4.0), distributed by arXiv under the Creative Commons Attribution 4.0 International licence, http://creativecommons.org/licenses/by/4.0/, as recorded in the arXiv licence metadata for this exact version and shown on the abstract page https://arxiv.org/abs/2601.04839v1. Full licence notice: \texttt{licenses/arxiv-2601.04839-CC-BY-4.0.txt}.\par\smallskip

\noindent\textit{Editorial scope.} Counted unit: the article's theorem Theorem11 (source0.tex lines 866--887). The article's complete original proof of it is delivered with it (source0.tex lines 889--1092, one \texttt{proof} environment of 204 lines). No mathematics is written, repaired or re-derived: the statement and the proof are the article's own lines, in the article's own notation, and no retained line is substituted or re-flowed.  Retained uncounted as the source-local context the proof needs: lines 447--461; lines 475--496; lines 502--504; lines 561--566; lines 581--585; lines 671--679; lines 717--727.  Nothing was omitted from any retained span, so the package carries no omission record.  No retained line contains a citation, so the wrapper carries no bibliography and no bibliography entry.  No retained line contains a figure environment, an image-inclusion command or drawing code, so no image is delivered and the package declares no asset dependency.  Editorial formatting only, in addition: an AMS \texttt{amsart} preamble that restates the article's own declarations and macros used by the retained lines, namely the environments \texttt{lemma}, \texttt{theorem} on separate counters, together with the standard packages those lines need.  The retained environments print in this document as Lemma 1, Theorem 1 in order of appearance, whereas the article prints the same items with the numbering of its own full document.  The complete native source of the article as distributed in the arXiv e-print for this exact version is bundled unchanged as \texttt{sources/192312/source0.tex}.
Let $\Omega$ be an open bounded Lipschitz domain in $\mathbb{R}^{m}$ ($m=2,3$) with polyhedral boundary $\partial \Omega$, and $T>0$.
For a given $\mathbf{F}\in (L^{2}((0,T);L^{2}(\Omega)))^{d\times d}$, we consider the following convection-diffusion problem:
\begin{equation}
	\begin{split}
		\left\{
		\begin{aligned}
			\partial_{t}\mathbf{U} - {\rm div}\big(\mathcal{D}  \nabla \mathbf{U}\big) + \bbeta \cdot \nabla \mathbf{U} + \mu \mathbf{U} &= \mathbf{F}  &&\text{in } (0,T] \times \Omega, \\
			\mathbf{U}(\bx, t) &= 0  &&\text{on } (0,T] \times \partial \Omega, \\
			\mathbf{U}(\cdot, 0) &= \mathbf{U}^{0} &&\text{in } \Omega,
		\end{aligned}
		\right.\label{CDR}
	\end{split}
\end{equation}
where $\mathcal{D}=(d_{ij})_{i,j=1}^{m}\in [L^{\infty}(\Omega)]^{m\times m}$ is symmetric and uniformly strictly positive definite a.e.~in $\Omega$, $\bbeta=( \beta_{i})_{i=1}^{m}\in L^{\infty}((0,T);W^{1,\infty}(\Omega))^{m}$,  and $\mu\in \mathbb{R}^{+}_{0}$, respectively, are the diffusion coefficient, the convective field, and the reaction coefficient. 
We will assume that $ {\rm div} \bbeta=0$ in $\Omega\times[0,T]$.

The standard weak formulation of  \eqref{CDR} reads as follows: Find
$\mathbf{U}\in L^\infty((0,T);(H^1_0(\Omega))^{d\times d,\textrm{sym}})\cap\left( H^1((0,T);H^{-1}(\Omega))\right)^{d\times d,\textrm{sym}}$ such that, 
for almost all $t\in (0,T)$ the following holds
\begin{equation}
	\begin{split}
		\left\{ \begin{array}{ll}  
			(\partial_{t}\mathbf{U},\mathbf{V})_{\Omega} + a(\mathbf{U},\mathbf{V}) = (\mathbf{F},\mathbf{V})_{\Omega} \hspace{1cm} \forall \mathbf{V}\in 	\left(H^{1}_{0}(\Omega)\right)^{d\times d, \textrm{sym}},\\
			\hspace{1.6cm} \mathbf{U}(\cdot,0) = \mathbf{U}^{0}.
		\end{array} \right.\label{eq82}
	\end{split}
\end{equation}
Here, the bilinear form $a(\cdot,\cdot)$ is defined by
\begin{equation}
	a(\mathbf{W},\mathbf{V}):=\left(\mathcal{D}\nabla \mathbf{W},\nabla \mathbf{V}\right)_{\Omega}+(\bbeta\cdot\nabla \mathbf{W},\mathbf{V})_{\Omega}+\mu( \mathbf{W},\mathbf{V})_{\Omega}.\label{eq81}
\end{equation}
In the above definition we have slightly abused the notation, as the convective term $\bbeta$ might depend on $t$, but unless the context requires it, we will always denote
this bilinear form by $a(\cdot,\cdot)$.  Since we have supposed that $\bbeta$ is solenoidal, then the bilinear form $a(\cdot,\cdot)$ is elliptic, in the sense that for $\mathbf{V}\in (H^1_0(\Omega))^{d\times d}$ it holds that $(\bbeta\cdot\nabla\mathbf{V},\mathbf{V})_\Omega=0$ and hence $a(\mathbf{V},\mathbf{V})=(\mathcal{D}\nabla\mathbf{V},\nabla\mathbf{V})_\Omega+\mu(\mathbf{V},\mathbf{V})_\Omega$.
More precisely, for each $t\in (0,T)$ the 
bilinear form  $a(\cdot,\cdot)$ induces the following ``energy'' norm in $	\left(H^{1}_{0}(\Omega)\right)^{d\times d}$
\begin{equation*}
	\|\mathbf{V} \|_{a}=\left( a(\mathbf{V},\mathbf{V})\right)^{\frac{1}{2}},\qquad t\in [0,T].
\end{equation*}

\begin{equation}
	\epsilon\leq \lambda_{k}(\bU(\bx,t))\leq \kappa, \hspace{0.5cm}k=1,\cdots,d,\hspace{0.5cm}\text{ for almost all}\ \  (\bx,t)\in \Omega\times(0,T),\label{eq14}
\end{equation}

\begin{equation}
	\|\mathbf{V}- \mathbf{I}_{h}\mathbf{V}\|_{0,K}
	+ h_{K} |\mathbf{V}-\mathbf{I}_{h}\mathbf{V}|_{1,K}
	\leq C h^{\ell +1}_{K}|\mathbf{V}|_{\ell +1,K}.
	\label{tensorlagrange}
\end{equation}

\begin{equation}
	\|\mathbf{V}\|^{2}_{0,\partial K} 
	\leq C \left(h_K^{-1}\|\mathbf{V}\|^2_{0,K}+h_K |\mathbf{V}|^{2}_{1,K}\right).
	\label{Tensortrace}
\end{equation}

With these notations, the finite element method used in this work reads as follows:
\begin{equation}
	\left\{ \begin{array}{ll}
		\text{For $1\leq n \leq N$, find $\mathbf{U}_h^n \in \mathbb{V}_{\calP}^{\epsilon,\kappa}$ such that}\vspace{.1cm}
		\\  
		(\delta \mathbf{U}_{h}^{n},\mathbf{V}_{h}-\mathbf{U}_{h}^{n})_{\Omega} + a_{J}(\mathbf{U}_{h}^{n},\mathbf{V}_{h}-\mathbf{U}_{h}^{n}) \geq (\mathbf{F}^{n},\mathbf{V}_{h}-\mathbf{U}_{h}^{n})_{\Omega} \hspace{1cm} \forall \mathbf{V}_{h}\in \mathbb{V}_{\calP}^{\epsilon,\kappa}, \vspace{.1cm}\\
		\hspace{4.8cm} \mathbf{U}_{h}^{0} = \mathbf{I}_{h}\mathbf{U}^{0} .
	\end{array} \right.\label{eq199}
\end{equation}

\begin{lemma}[Discrete Gr\"onwall lemma]\label{13}
	Let $k$, $B$, $a_{n}$, $b_{n}$, $c_{n}$, $\gamma_{n}$, $n=0,\ldots,\nu$, be non-negative numbers such that
	\begin{align*}
		a_{\nu}+k\sum_{n=0}^{\nu}b_{n}\leq k\sum_{n=0}^{\nu}\gamma_{n}a_{n}+k\sum_{n=0}^{\nu}c_{n}+B
		\hspace{0.5cm}{\rm for}\hspace{0.3cm}\nu\geq 0.
	\end{align*}
	Suppose $k\gamma_{n}\leq 1$ for every $n$, and set $\sigma_{n}=(1-k\gamma_{n})^{-1}$. Then
	\begin{align}
		a_{\nu}+k\sum_{n=0}^{\nu}b_{n}\leq \exp\left(k\sum_{n=0}^{\nu}\sigma_{n}\gamma_{n}\right)\left(k\sum_{n=0}^{\nu}c_{n}+B\right).\label{inequality11}
	\end{align}
\end{lemma}

\begin{theorem}[A priori error estimate]
  \label{Theorem11}
	Let $k\ge 1$, let $\bU^{0}\in \left(H^{k+1}(\Omega)\right)^{d\times d}$, and let $\bU$ be the solution of \eqref{CDR} satisfying
	\[
	\bU\in  L^{\infty}((0,T);(H^{k+1}(\Omega))^{d\times d})\cap (H^{1}(0,T);(H^{k+1}(\Omega))^{d\times d})\cap H^{2}((0,T);(L^{2}(\Omega))^{d\times d}),
	\]
	with $\bU(\cdot,t)\in (H^{1}_{0}(\Omega))^{d\times d,\textrm{sym}}$ for almost all $t\in [0,T]$.
	Assume in addition that \eqref{eq14} holds for $\bU(t)$ for almost all $t\in[0,T]$, so that $\mathbf I_h\bU^n\in\mathbb V_{\calP}^{\epsilon,\kappa}$ for each $n$.
	Let $\bU_{h}^{n}\in \mathbb{V}_{\calP}$ be the solution of \eqref{eq199} 
	at the time step $n$.
	Defining $\bE^{n}=\bU^{n}_{h}-\bU^{n}$, there exists a constant $C>0$,  independent of $h$ and $\Delta t$ (and depending only on $k$ and mesh shape-regularity), such that
\begin{equation}
	\begin{aligned}
		\Big( \|\bE^N \|_{0,\Omega}^2&+\Delta t\sum_{n=1}^{N}\|\bE^n \|_{a_{J}}^2 \Big)^{\frac{1}{2}}
		\leq C\Bigg[ \Delta t \left\|	\partial_{t}^{2}\mathbf{U} \right\|_{L^{2}((0,T);L^{2}(\Omega))} \\
		&\qquad\qquad+ h^{k}\Big(\big(T\|\bbeta\|_{0,\infty,\Omega}^{2}+\|\mathcal{D}^{\frac{1}{2}}\|_{0,\infty,\Omega}^{2}+\gamma h\|\bbeta\|_{0,\infty,\Omega}+\mu h^{2}\big)^{\frac{1}{2}}
		\left\|	\mathbf{U} \right\|_{L^{2}((0,T);H^{k+1}(\Omega))}\\
			&\qquad\qquad\qquad\qquad\qquad\qquad
			+h\Big(\left\|\partial_{t}\mathbf{U} \right\|_{L^{2}((0,T);H^{k+1}(\Omega))}+ \max\limits_{1\leq n\leq N}|\bU^{n} |_{k+1,\Omega}\Big)\Big)\Bigg].
	\end{aligned}
\end{equation}
\end{theorem}

\begin{proof}
	Let $\mathbf{U}^n = \mathbf{U}(t_n)$ denote the exact solution at time $t_n$.
	Using the test function $\mathbf{I}_h \mathbf{U}^n\in \mathbf{V}_{\calP}^{\epsilon,\kappa}$ in \eqref{eq199} gives
	\begin{align}
		\begin{split}\label{inequalii}
			\left(\delta\mathbf{U}_h^n, \mathbf{I}_h\mathbf{U}^n - \mathbf{U}_h^n\right)_{\Omega}
			+  a_J(\mathbf{U}_h^{n}, \mathbf{I}_h\mathbf{U}^n - \mathbf{U}_h^n)
			\geq (\mathbf{F}^{n}, \mathbf{I}_h\mathbf{U}^n - \mathbf{U}_h^n)_{\Omega}.
		\end{split}
	\end{align}
	Also, setting $\mathbf{V}_h = \mathbf{I}_h \mathbf{U}^n - \mathbf{U}_h^n$ in the continuous problem \eqref{eq82} yields
	\begin{align}
		\begin{split}\label{exact12}
			\left(\partial_t\mathbf{U}(t_{n}), \mathbf{I}_h \mathbf{U}^n - \mathbf{U}_h^n\right)_{\Omega}
			+ a(\mathbf{U}(t_{n}), \mathbf{I}_h \mathbf{U}^n - \mathbf{U}_h^n)
			=(\mathbf{F}(t_{n}), \mathbf{I}_h \mathbf{U}^n - \mathbf{U}_h^n)_{\Omega}.
		\end{split}
	\end{align}

	Now, we decompose $\bE^{n}=\bU^{n}_{h}-\bU^{n}$ as
\begin{align}
	\bE^{n}=\left(\bU^{n}_{h}-\mathbf{I}_h \mathbf{U}^n\right)+\left(\mathbf{I}_{h}\bU^{n}-\bU^{n}\right)=:\bE^{n}_{h}+\eta^{n}_{h}.\label{error1}
\end{align}
Since $k\ge 1$ and $\bU^{n}\in (H^{k+1}(\Omega))^{d\times d}\subset (H^{2}(\Omega))^{d\times d}$, we have $\llbracket \nabla \bU^{n}\rrbracket=\mathbf 0$ on every interior facet and hence
\[
	J(\bU^{n},\bV_h)=0\qquad \forall \bV_h\in\mathbb V_{\calP}.
\]
Consequently,
\begin{align}
	J(\bU^{n}_{h}, \bV_{h})
	=J(\bU^{n}_{h}- \bU^{n}, \bV_{h})
	=J( \bE^{n}_{h}, \bV_{h})+J(\eta^{n}_{h}, \bV_{h}).\label{CIP12}
\end{align}
Adding $J(\bU^n,\mathbf I_h\bU^n-\bU_h^n)=0$ to \eqref{exact12}, we may equivalently replace $a(\bU^n,\cdot)$ by $a_J(\bU^n,\cdot)$ in \eqref{exact12}.
Therefore, subtracting \eqref{exact12} from \eqref{inequalii} gives
\begin{align}\label{subtraction}
	\Big(\delta\mathbf{U}_h^n - \partial_t\mathbf{U}(t_{n}), \mathbf{I}_h \mathbf{U}^n - \mathbf{U}_h^n\Big)_{\Omega}
	+ a_J(\bU_h^{n}-\bU(t_{n}), \mathbf{I}_h \mathbf{U}^n - \mathbf{U}_h^n)
	\geq 0.
\end{align}

The first term in \eqref{subtraction} may be rewritten as
\begin{align}
	\delta \bU^{n}_{h}-\partial_{t} \bU(t_{n})
	&=\big(\delta \bU^{n}_{h}- \delta(\mathbf{I}_{h}\bU^{n})\big)
	+\big(\delta\bU^{n}-\partial_{t} \bU(t_{n})\big)
	+\big(\delta(\mathbf{I}_{h}\bU^{n})-\delta\bU^{n}\big)\nonumber\\
	&= \delta \bE^{n}_{h}+\big(\delta\bU^{n}-\partial_{t} \bU(t_{n})\big)+\delta \eta^{n}_{h}. \label{time12}
\end{align}
Using \eqref{time12} and $\mathbf I_h\bU^n-\bU_h^n=-\bE_h^n$, \eqref{subtraction} becomes
\begin{align}\label{Eq12367443}
	\big(\delta \bE^{n}_{h}, -\bE^{n}_{h}\big)_{\Omega}+a_{J}(\bE^{n}_{h},-\bE^{n}_{h})
	\geq
	(\partial_{t} \bU(t_{n})-\delta \bU^{n},-\bE^{n}_{h})_{\Omega}
	-(\delta \eta^{n}_{h}, -\bE^{n}_{h})_{\Omega}
	-a_{J}( \eta^{n}_{h}, -\bE^{n}_{h}).
\end{align}
Multiplying \eqref{Eq12367443} by $-1$ gives the equivalent form
\begin{align}\label{Eq12367443b}
	\big(\delta \bE^{n}_{h}, \bE^{n}_{h}\big)_{\Omega}+a_{J}(\bE^{n}_{h},\bE^{n}_{h})
	\leq
	(\delta \bU^{n}-\partial_{t} \bU(t_{n}),\bE^{n}_{h})_{\Omega}
	+(\delta \eta^{n}_{h}, \bE^{n}_{h})_{\Omega}
	+a_{J}( \eta^{n}_{h}, \bE^{n}_{h}).
\end{align}

Using the identity $2(a^{2}-ab)=a^{2}-b^{2}+(a-b)^{2}$ in $(\delta\bE_h^n,\bE_h^n)_\Omega$ and the Cauchy--Schwarz inequality on the right-hand side of \eqref{Eq12367443b}, we obtain
\begin{align*}
	&\frac{1}{2\Delta t}\Big(\|\bE_h^n \|_{0,\Omega}^2-\| \bE_h^{n-1}\|_{0,\Omega}^2+\|\bE_h^n - \bE_h^{n-1}\|_{0,\Omega}^2\Big)
	+ \|\bE_h^n\|_{a_J}^2\\
	\leq&\, \Big(\|\delta \bU^{n}-\partial_{t}\bU(t_n)\|_{0,\Omega}+\|\delta\eta_h^n\|_{0,\Omega}
	+\|\bbeta\|_{0,\infty,\Omega} |\eta_h^n|_{1,\Omega}\Big)\, \|\bE_h^n\|_{0,\Omega} + \|\eta_h^n\|_{a_J} \|\bE_h^n\|_{a_J}.
\end{align*}
Applying Young's inequality yields
\begin{align*}
	&\|\bE_h^n \|_{0,\Omega}^2-\| \bE_h^{n-1}\|_{0,\Omega}^2+\|\bE_h^n - \bE_h^{n-1}\|_{0,\Omega}^2
	+ \Delta t \|\bE_h^n\|_{a_J}^{2}\\
	\leq&\,  2\Delta t\Big(T(\|\delta \bU^{n}-\partial_{t}\bU(t_n)\|_{0,\Omega}+\|\delta\eta_h^n\|_{0,\Omega}
	+\|\bbeta\|_{0,\infty,\Omega} |\eta_h^n|_{1,\Omega})^2+\frac{1}{T}\|\bE_h^n\|_{0,\Omega}^2\Big)
	+\Delta t \| \eta_h^n\|_{a_J}^{2}.
\end{align*}

Summing from $n=1$ to $n=N$ and using $\bE_h^{0}=\bU_h^0-\mathbf I_h\bU^0=\mathbf 0$ gives
\begin{equation}
	\begin{aligned}\label{error432}
	&\|\bE_h^N \|_{0,\Omega}^2
	+\sum_{n=1}^{N}\|\bE_h^n - \bE_h^{n-1}\|_{0,\Omega}^2
	+ \Delta t \sum_{n=1}^{N}\|\bE_h^n \|_{a_{J}}^{2}\\
	\leq&\, C\Delta t\sum_{n=1}^{N}\bigg(
	T\Big(\|\delta \bU^{n}-\partial_{t}\bU(t_n)\|_{0,\Omega}^{2}
	+\|\delta\eta_h^n\|_{0,\Omega}^{2}
	+\|\bbeta\|_{0,\infty,\Omega}^{2} |\eta_h^n|_{1,\Omega}^{2}\Big)
	+\|\eta_h^n\|_{a_J}^{2}
	+\frac{1}{T}\|\bE_h^n\|_{0,\Omega}^{2}\bigg).
	\end{aligned}
\end{equation}

We now apply Gr\"{o}nwall's Lemma~\ref{13} with the choices
\[
	k=\Delta t,\qquad
	a_n=\|\bE_h^n\|_{0,\Omega}^2,\qquad
	\gamma_n=\frac{1}{T},\qquad
	\sigma_n=\Big(1-\frac{\Delta t}{T}\Big)^{-1},
\]
\[
	b_n=\frac{1}{\Delta t}\|\bE_h^n-\bE_h^{n-1}\|_{0,\Omega}^2+\|\bE_h^n\|_{a_J}^{2},
	\]
\[
	c_n=C\Big(T(\|\delta \bU^{n}-\partial_{t}\bU(t_n)\|_{0,\Omega}^{2}
	+\|\delta\eta_h^n\|_{0,\Omega}^{2}
	+\|\bbeta\|_{0,\infty,\Omega}^{2} |\eta_h^n|_{1,\Omega}^{2})
	+\|\eta_h^n\|_{a_J}^{2}\Big),
\]
and $B=0$. Since $\Delta t=T/N$ and $N\ge 2$, we have $k\gamma_n=\Delta t/T=1/N<1$ and hence
\[
	\exp\Big(k\sum_{n=1}^{N}\sigma_n\gamma_n\Big)
	=\exp\Big(\frac{N}{N-1}\Big)\le e^{2}.
\]
Therefore,
\begin{align}
	&\|\bE_h^N \|_{0,\Omega}^2+\Delta t \sum_{n=1}^{N}\|\bE_h^n \|_{a_{J}}^{2} \nonumber\\
	\leq&\,  C e^{2} \Delta t\sum_{n=1}^{N}\Big(
	T\|\delta \bU^{n}-\partial_{t}\bU(t_n)\|_{0,\Omega}^{2}
	+T\|\delta\eta_h^n\|_{0,\Omega}^{2}
	+T\|\bbeta\|_{0,\infty,\Omega}^{2} |\eta_h^n|_{1,\Omega}^{2}
	+\|\eta_h^n\|_{a_J}^{2}\Big).\label{53last}
\end{align}

For the time truncation term, Taylor's theorem gives, for each $n$,
\[
	\partial_t \bU(t_n)-\delta \bU^n
	=\frac{1}{\Delta t}\int_{t_{n-1}}^{t_n}(t_n-s) \partial_t^2\bU(s) {\rm d}s,
\]
and hence, by Cauchy--Schwarz,
\begin{equation}
	\sum_{n=1}^{N} \left\|\partial_t\mathbf{U}(t_{n})-\delta \mathbf{U}^n\right\|_{0,\Omega}^2
	\leq C\Delta t\int_{0}^{T} \left\|\partial_{t}^{2}\mathbf{U}(t)\right\|_{0,\Omega}^2 {\rm d}t.  \label{time_trunc}
\end{equation}

Next, using the tensor Lagrange approximation \eqref{tensorlagrange}, we have
\begin{align}
	|\eta^{n}_{h}|_{1,\Omega}^{2}\leq Ch^{2k}|\bU^{n} |_{k+1,\Omega}^{2},
	\qquad
	\|\eta^n_h\|_{0,\Omega}^2\le C  h^{2k+2}|\bU^n|_{k+1,\Omega}^{2}.\label{inequal32}
\end{align}
Moreover,
\begin{equation}
	\begin{aligned}
		\sum_{n=1}^{N}\|\delta \eta_{h}^{n}\|_{0,\Omega}^{2}
		&=\sum_{n=1}^{N}\Big\|\delta (\mathbf{I}_h\bU^{n})-\delta(\bU^{n})\Big\|_{0,\Omega}^{2}
		\leq Ch^{2k+2}\sum_{n=1}^{N}|\delta \mathbf{U}^{n} |_{k+1,\Omega}^{2}\\
		&\leq C h^{2k+2}\sum_{n=1}^{N}\left|\frac{1}{\Delta t}\int_{t_{n-1}}^{t_{n}}\partial_{t} \mathbf{U}(t) {\rm d}t\right|_{k+1,\Omega}^{2}
		\leq C\frac{h^{2k+2}}{\Delta t}\int_{0}^{T}|\partial_{t}\mathbf{U}(t)|_{k+1,\Omega}^{2}{\rm d}t.
	\end{aligned}
\end{equation}

Finally, for the stabilisation part we use the discrete trace inequality \eqref{Tensortrace} (applied componentwise) together with standard scaling and the interpolation estimates \eqref{tensorlagrange} to obtain
\begin{align}
	J(\eta_{h}^{n},\eta_{h}^{n})
	= \gamma \sum_{F\in\calF_I^{}}\int_F\|\bbeta \|_{0,\infty,F}  h_{F}^{2} \llbracket\nabla \eta_{h}^{n} \rrbracket :\llbracket\nabla \eta_{h}^{n} \rrbracket  {\rm d}s
	\leq C\gamma h^{2k+1}\|\bbeta\|_{0,\infty,\Omega} |\bU^{n} |_{k+1,\Omega}^{2}.\label{JJJ}
\end{align}
Using \eqref{inequal32} and \eqref{JJJ}, we obtain
\begin{equation}
	\begin{aligned}
		T\|\bbeta\|_{0,\infty,\Omega}^{2} |\eta^{n}_{h}|_{1,\Omega}^{2}
		+\|\eta^{n}_{h}\|_{a_J}^{2}
		&\leq Ch^{2k}\Big(\big(T\|\bbeta\|_{0,\infty,\Omega}^{2}
		+\|\mathcal{D}^{\frac{1}{2}}\|_{0,\infty,\Omega}^{2}
		+\gamma h\|\bbeta\|_{0,\infty,\Omega}
		+\mu h^{2}\big) |\mathbf{U}^{n}|_{k+1,\Omega}^{2}\Big).
	\end{aligned}
\end{equation}
Summing in $n$ and using Cauchy--Schwarz yields
\begin{align}\label{eq:summed_interp_bound}
		&\sum_{n=1}^{N}\Big(T\|\bbeta\|_{0,\infty,\Omega}^{2} |\eta^{n}_{h}|_{1,\Omega}^{2}
		+\|\eta^{n}_{h}\|_{a_J}^{2}\Big) \nonumber\\
		\leq&\,  \frac{Ch^{2k}}{\Delta t}\big(T\|\bbeta\|_{0,\infty,\Omega}^{2}
		+\|\mathcal{D}^{\frac{1}{2}}\|_{0,\infty,\Omega}^{2}
		+\gamma h\|\bbeta\|_{0,\infty,\Omega}
		+\mu h^{2}\big)\int_{0}^{T}|\mathbf{U}(t)|_{k+1,\Omega}^{2}{\rm d}t.
\end{align}

Inserting \eqref{time_trunc}, \eqref{eq:summed_interp_bound} and the bound for $\sum_n\|\delta\eta_h^n\|_{0,\Omega}^2$ into \eqref{53last}, and rearranging, gives
\begin{equation}
	\begin{aligned}
		\max_{1\leq n\leq N}\|\bE_h^n \|_{0,\Omega}^2+\Delta t\sum_{n=1}^{N}\|\bE_h^n \|_{a_{J}}^2
		\leq C\Bigg[
		\Delta t^{2}\int_{0}^{T} \left\|\partial_{t}^{2}\mathbf{U}\right\|_{0,\Omega}^2 {\rm d}t\\
		\qquad\qquad
		+h^{2k}\Big(\big(T\|\bbeta\|_{0,\infty,\Omega}^{2}+\|\mathcal{D}^{\frac{1}{2}}\|_{0,\infty,\Omega}^{2}+\gamma h\|\bbeta\|_{0,\infty,\Omega}+\mu h^{2}\big)\int_{0}^{T}|\mathbf{U}(t)|_{k+1,\Omega}^{2}{\rm d}t\\
		\qquad\qquad\qquad\qquad
		+h^{2}\int_{0}^{T}|\partial_{t}\mathbf{U}(t)|_{k+1,\Omega}^{2}{\rm d}t\Big)\Bigg].
	\end{aligned}
\end{equation}

Finally, using the triangle inequality,
\[
	\|\bE^n\|_{0,\Omega} \le \|\bE_h^n\|_{0,\Omega}+\|\eta_h^n\|_{0,\Omega},
	\qquad
	\|\bE^n\|_{a_J} \le \|\bE_h^n\|_{a_J}+\|\eta_h^n\|_{a_J},
\]
together with \eqref{inequal32}, yields the stated estimate after taking square roots.
\end{proof}

\end{document}
